\subsection{Уравнения линии насыщения и границы B23}

	Линия насыщения (регион~4) является одномерным множеством состояний,
в которых жидкая и паровая фазы находятся в равновесии.
Физически это означает,
что по обе стороны линии располагаются однофазные области,
а внутри купола насыщения существует двухфазная смесь.
Поэтому для входной пары $(p, T)$
оба параметра на линии насыщения уже не являются независимыми:
задание одного из них однозначно определяет второй.
Формально линия насыщения соответствует условию фазового равновесия,
при котором удельные энергии Гиббса жидкой и паровой фаз совпадают.

\textbf{Тройная и критическая точки.}
	Нижней опорной точкой термодинамической шкалы
	служит тройная точка воды
	с параметрами $T_{t}=273{,}16$~К и $p_{t}=611{,}657$~Па,
	в которой сосуществуют твёрдая, жидкая и паровая фазы.
	Именно к этой точке привязана нормировка региона~1,
где удельные внутренняя энергия и энтропия
насыщенной жидкости принимаются равными нулю.
	Верхней точкой линии насыщения
	служит критическая точка
	$T_{c}=647{,}096$~К,
	$p_{c}=22{,}064$~МПа,
	$\rho_{c}=322$~кг/м\textsuperscript{3}.
В критической точке различие между насыщенной жидкостью
и насыщенным паром исчезает,
а разность их удельных объёмов стремится к нулю.
При этом выполняются условия
\begin{equation}
\left(\frac{\partial p}{\partial \rho}\right)_{T} = 0,
\qquad
\left(\frac{\partial^{2} p}{\partial \rho^{2}}\right)_{T} = 0,
\end{equation}
что и делает околокритическую область
особенно чувствительной к численным ошибкам.

\textbf{Аналитическое задание линии насыщения.}
Для региона~4 стандарт IF97 использует единое неявное полиномиальное соотношение
~\cite{iapws_if97_2012}:
\begin{equation}
\beta^{2}\vartheta^{2}
+ n_{1}\beta^{2}\vartheta
+ n_{2}\beta^{2}
+ n_{3}\beta\vartheta^{2}
+ n_{4}\beta\vartheta
+ n_{5}\beta
+ n_{6}\vartheta^{2}
+ n_{7}\vartheta
+ n_{8}
= 0,
\end{equation}
\begin{gostwhere}
  \item $\beta$~--- безразмерная переменная давления, $\beta = \left(\frac{p_{\mathrm{sat}}}{p^{*}}\right)^{1/4}$;
  \item $\vartheta$~--- вспомогательная безразмерная переменная температуры, $\vartheta = \frac{T}{T^{*}} + \frac{n_{9}}{\frac{T}{T^{*}} - n_{10}}$;
  \item $p^{*}$~--- нормировочное давление, $p^{*}=1~\mbox{МПа}$;
  \item $T^{*}$~--- нормировочная температура, $T^{*}=1~\mbox{К}$.
\end{gostwhere}

Все коэффициенты $n_{1},\dots,n_{10}$
задаются стандартом в табличной форме.

Из этого полинома выводится прямое уравнение насыщения
для зависимости $p_{\mathrm{sat}}(T)$:
\begin{equation}
p_{\mathrm{sat}}(T) = p^{*}
\left[
\frac{2C}{-B + \sqrt{B^{2} - 4AC}}
\right]^{4},
\end{equation}
\begin{gostwhere}
  \item $A$~--- вспомогательный коэффициент, $A = \vartheta^{2} + n_{1}\vartheta + n_{2}$;
  \item $B$~--- вспомогательный коэффициент, $B = n_{3}\vartheta^{2} + n_{4}\vartheta + n_{5}$;
  \item $C$~--- вспомогательный коэффициент, $C = n_{6}\vartheta^{2} + n_{7}\vartheta + n_{8}$.
\end{gostwhere}

Аналогично получается и обратное уравнение
для зависимости $T_{\mathrm{sat}}(p)$:
\begin{equation}
T_{\mathrm{sat}}(p) = \frac{T^{*}}{2}
\left[
n_{10} + D - \sqrt{(n_{10}+D)^{2} - 4(n_{9} + n_{10}D)}
\right],
\end{equation}
\begin{gostwhere}
  \item $D$~--- вспомогательная величина, $D = \frac{2G}{-F - \sqrt{F^{2} - 4EG}}$;
  \item $E$~--- вспомогательный коэффициент, $E = \beta^{2} + n_{3}\beta + n_{6}$;
  \item $F$~--- вспомогательный коэффициент, $F = n_{1}\beta^{2} + n_{4}\beta + n_{7}$;
  \item $G$~--- вспомогательный коэффициент, $G = n_{2}\beta^{2} + n_{5}\beta + n_{8}$.
\end{gostwhere}

Таким образом, линия насыщения
задаётся аналитически,
а переход между постановками $(T)$ и $(p)$
не требует итерационного решения.

\textbf{Почему регион~4 является численно сложной областью.}
Внутри двухфазной области свойства среды
не задаются одним фундаментальным уравнением состояния.
Вместо этого требуется вычислить свойства насыщенной жидкости
и насыщенного пара,
а затем выполнить смешение по степени сухости $x$.
Кроме того, на линии насыщения и вблизи критической точки
производные свойств ведут себя жёстко:
изобарная теплоёмкость быстро возрастает,
скорость звука демонстрирует выраженный провал,
а величина $v''-v'$ убывает,
что ухудшает обусловленность обратных задач.

\textbf{Как эта сложность обрабатывается в модели.}
В проекте проблемы региона~4 решаются комбинированно.
Сама линия насыщения вычисляется по аналитическим формулам
$p_{\mathrm{sat}}(T)$ и $T_{\mathrm{sat}}(p)$.
Свойства насыщенной жидкости и насыщенного пара
ниже $623{,}15$~К рассчитываются по формулам регионов~1 и~2,
а в околокритической части линии насыщения ---
через формулировку региона~3.
Для смеси используются явные правила смешения по $x$,
а при вычислениях на границах вводятся проверки диапазонов,
контроль принадлежности куполу насыщения
и отдельная обработка плохо обусловленного поведения
вблизи критической точки.

Граница B23 отделяет регион~2 от региона~3.
Она определяет момент перехода
от формулировки через потенциал Гиббса
к формулировке через потенциал Гельмгольца.
В безразмерной записи эта граница задаётся
квадратичной аппроксимацией вида
\begin{equation}
\pi = n_{1} + n_{2}\theta + n_{3}\theta^{2},
\end{equation}
\begin{gostwhere}
  \item $\pi$~--- безразмерное давление, $\pi = p/p^{*}$;
  \item $\theta$~--- безразмерная температура, $\theta = T/T^{*}$~\cite{iapws_if97_2012}.
\end{gostwhere}

Корректная обработка B23 обязательна
для маршрутизации в постановках $(p, T)$, $(p, h)$ и $(p, s)$.
